\section{Final comparison of the counts}\label{balance}
\begin{figure}[htbp]
\centering
\begin{tikzpicture}[x=1cm,y=1cm,>=Stealth]
 \draw[->,line width=.6pt] (0,0) -- (13.5,0);
 \foreach \x/\lab in {0/{k=n},4.4/{k=M=\lceil\delta n\rceil},9.5/{k=K}}{
  \draw[line width=.6pt] (\x,.12) -- (\x,-.12);
  \node[above,font=\small] at (\x,.14) {$\lab$};
 }
 \foreach \x in {4.4,9.5}{\draw[black!35,line width=.3pt] (\x,-.12) -- (\x,-3.25);}
 \node[font=\small,align=center,text width=4.1cm] at (2.2,-.62)
  {Macroscopic range\\ (Section~\ref{macro})};
 \node[font=\small,align=center,text width=4.8cm] at (6.95,-.62)
  {Mesoscopic range\\ (Section~\ref{meso})};
 \node[font=\small,align=center,text width=3.6cm] at (11.5,-.62)
  {Last board\\ (Section~\ref{terminal})};
 \foreach \y/\lab in {-1.55/Regularity,-2.25/Exit,-2.95/Cost}{
  \node[font=\footnotesize,anchor=east] at (-.25,\y) {\lab};
 }
 \node[font=\footnotesize] at (2.2,-1.55) {$\Delta<\varrho_{\rm mac}$};
 \node[font=\footnotesize] at (6.95,-1.55) {inner class};
 \node[font=\footnotesize] at (11.5,-1.55) {event $F$};
 \node[font=\footnotesize] at (2.2,-2.25) {$\nu(E_{\rm mac})\le Ce^{-cn}$};
 \node[font=\footnotesize] at (6.95,-2.25) {$\nu(E_j)\le Ce^{-c\,2^jK}$};
 \node[font=\footnotesize] at (11.5,-2.25) {$\nu(F\cap D_K)\le d_{n,K}$};
 \node[font=\footnotesize] at (2.2,-2.95) {$g^+\le C_\delta n^{2/3}$};
 \node[font=\footnotesize] at (6.95,-2.95)
  {$g^+\le C_0+2\beta(\varepsilon_{\rm ext})2^jK$};
 \node[font=\footnotesize] at (11.5,-2.95) {$|g|\le\theta_K$ on $F\cap D_K^c$};
\end{tikzpicture}
\caption{The three ranges of a history, from left to right as the size $k$ decreases. The first
row gives the regularity condition of Definition~\ref{def:regular}. The second gives the bound on
the exit probability under $\nu$, with the events $E_{\rm mac}$ and $E_j$ of this section
(Corollary~\ref{macro:cambio}, Lemma~\ref{meso:entrada}, and
Proposition~\ref{prop:salida-banda-meso}) and, on the last board, the bound of
Lemma~\ref{terminal-geometria}. The third gives the bound on the cost at the first-exit board
(Proposition~\ref{coste-macro} and Corollary~\ref{coste-debil}) and the bound of
Proposition~\ref{terminal-error}. The constants $C$ and $c$ in the macroscopic range depend on
$\delta$ and $\varrho_{\rm mac}$. Corollary~\ref{hist:consistencia} bounds $A_\tau$ along the whole
history.}\label{fig:tramos}
\end{figure}

\begin{proof}[Proof of Theorem~\ref{princ}]
Fix the class constants, $a_{\rm mes}$, $\varrho_{\rm mes}$ and the tolerances, $\varepsilon$,
$\delta$, $K_0$, $\varrho_{\rm mac}$ and $a_{\rm mac}$ in the order given in
Section~\ref{meso:constantes}. Let $K\ge K_0$ be a fixed integer and let $\tau$ be the stopping
time in Definition~\ref{def:regular}. This stopping time satisfies the hypotheses of
Section~\ref{historias}. Indeed, every parent before $\tau$ is regular and therefore has a Bethe
point with fixed bounds: by Proposition~\ref{macro:maximo} and Corollary~\ref{macro:cambio} in the
macroscopic range, by Proposition~\ref{meso:reconstruccion} at the cut-off, and by definition in
the mesoscopic range. By item (I1) of Section~\ref{leyes}, each of its children has a Bethe point
and a finite predictor. These same bounds are the hypotheses of
Corollary~\ref{hist:consistencia}. Let $E_{\rm mac}$ be the event that the first-exit board has
size ${k\ge M}$; this includes failure of entry at the cut-off $k=M$. For $j\ge0$, let $E_j$ be
the event that the history has regular entry at $M$ and its first-exit board has size $k$ with
${2^jK\le k<2^{j+1}K}$; this is the event in Proposition~\ref{prop:salida-banda-meso}. If $K<M$,
an exit on the board of size $K$ belongs to $E_0$. These events and the permanence event $F$
partition the stopped histories.

On $E_{\rm mac}$ and on each $E_j$, the first-exit board $B_\tau$ inherits the outer bounds from
its regular parent; on $E_j$ it also inherits diagonal mass at most $\varepsilon_{\rm ext}$
(Proposition~\ref{meso:reconstruccion}). On $E_{\rm mac}$ its size is $k\ge M\ge\delta n$, so
$N=2n-1=O_\delta(k)$ and Proposition~\ref{coste-macro} bounds its exit cost by
$\Lambda_{\rm mac}=C_\delta n^{2/3}$. On $E_j$, writing $k$ for its size,
Corollary~\ref{coste-debil} gives
\[
 g(B_\tau)\le C_0+\beta(\varepsilon_{\rm ext})k
           \le C_0+2\beta(\varepsilon_{\rm ext})2^jK=:\Lambda_j .
\]
Both bounds hold on histories of positive mass under $\mu$. Let $\Lambda$ be the function of the
stopped history equal to $\Lambda_{\rm mac}$ on $E_{\rm mac}$, to $\Lambda_j$ on $E_j$, and to zero
on $F$. Thus $g(B_\tau)\mathbf1_{F^c}\le\Lambda$ on these histories.

We apply the Gibbs variational inequality~\cite[Proposition~1.4.2]{de} (see
also~\cite[Ch.~2]{ct}) to $\Lambda$ and to the laws of the stopped
histories. On this finite space the inequality follows directly from Jensen's inequality by
restricting the sum under $\nu$ to the support of $\mu$:
\[
 \E_\mu\Lambda-D(\mu^\tau\Vert\nu^\tau)\le\log\E_\nu e^\Lambda
 =\log\!\left(1+\nu(E_{\rm mac})(e^{\Lambda_{\rm mac}}-1)
              +\sum_{j\ge0}\nu(E_j)(e^{\Lambda_j}-1)\right).
\]
If the first exit occurs at $M$ with $\Delta<\varrho_{\rm mac}$, Lemma~\ref{meso:entrada} places
the observables within their entry tolerances, so $D_{\max}(M)\ge\alpha_{\rm geo} M$. Thus
\[
 E_{\rm mac}\subseteq\mathcal H_M\cup\bigl(G_M^-\cap\{D_{\max}(M)\ge\alpha_{\rm geo} M\}\bigr),
\]
so Corollary~\ref{macro:cambio} and that lemma bound the macroscopic summand by
${C_1\exp(-c_1n+C_\delta n^{2/3})}$, with $C_1,c_1>0$ depending on $\delta$ and $\varrho_{\rm mac}$.
Proposition~\ref{prop:salida-banda-meso} gives ${\nu(E_j)\le C\exp(-c\,2^jK)}$ if $2^jK\le M$;
otherwise $E_j$ is empty. Since $c>2\beta(\varepsilon_{\rm ext})$, for
$0<c'<c-2\beta(\varepsilon_{\rm ext})$ the sum of $\nu(E_j)(e^{\Lambda_j}-1)$ is $O(e^{-c'K})$ for
large $K$. Corollary~\ref{hist:consistencia}, \eqref{super}, and the terminal estimate
under $\mu$ in Proposition~\ref{terminal-error} give
\[
 \limsup_{n\to\infty}g(B_0)\le C/K+\theta_K+Ce^{-c'K},\qquad\theta_K\to0 .
\]
Indeed, for fixed $K$ the exceptional cost on the last board is at most
${(L_K^{\max}+\log K!)\,CK^4/n}$, which is ${O(K^5\log(K+1)/n)}$ and tends to zero before we take
the limit in $K$.

For the reverse inequality we use $G=F\cap\{Q(B^{(K)})>0\}$. By the geometric estimate under $\nu$
in Section~\ref{terminal}, $s=\nu(G)\ge1-Ce^{-cK}-o_n(1)-d_{n,K}$, where $d_{n,K}\to0$ for fixed
$K$. For large $K$ this lower limit is positive. The negative part of the exceptional cost on the
last board has expectation at most $L_K^{\max}d_{n,K}/s$ under $\nu(\cdot\mid G)$. On the remaining
histories $\theta_K$ suffices. Taking the lower limit in \eqref{infer} and using
$\log(1-z)\ge-2z$ for $0\le z\le1/2$ gives
\[
 \liminf_{n\to\infty}g(B_0)\ge-C/K-\theta_K-Ce^{-cK}.
\]
Now let $K$ tend to infinity in both inequalities. On the initial board $B_0$, the global maximizer
in Proposition~\ref{macro:maximo} satisfies $p_{\max}<1/4$ for large $n$ and hence coincides with
the canonical maximizer of~\cite{pred}. The relation between the cofactor and the pseudodeterminant
in Section~\ref{leyes} gives $L(B_0)=\log Q_{\rm pred}(n)$. Thus
$g(B_0)=\log Q(n)-\log Q_{\rm pred}(n)\to0$. Applying Theorem~1.1 of~\cite{pred} to the second
logarithm, which is item (I5), identifies the constant in the statement.
\end{proof}
